\documentclass[10pt,twocolumn]{article}

\usepackage[T1]{fontenc}
\usepackage[utf8]{inputenc}
\usepackage{lmodern}
\usepackage{microtype}

\usepackage[letterpaper,margin=1in,columnsep=0.3in]{geometry}
\setlength{\parindent}{1em}
\setlength{\parskip}{0pt}

\usepackage{amsmath,amssymb,amsthm}
\usepackage{graphicx}
\usepackage{booktabs}
\usepackage{array}
\usepackage{tikz}
\usepackage{tikz-3dplot}
\usetikzlibrary{calc}

\definecolor{tfteal}{HTML}{00D5BE}
\definecolor{tfrose}{HTML}{C8467A}
\definecolor{tfink}{HTML}{1F2937}

\usepackage[numbers,sort&compress]{natbib}
\usepackage[hidelinks]{hyperref}
\hypersetup{
  pdftitle={Exact predicates, implicit graphs, and aggregated uncertainty in mesh CSG},
  pdfauthor={\v{Z}iga Sajovic},
}

\title{Exact predicates, implicit graphs,\\
       and aggregated uncertainty in mesh CSG}
\author{%
  \v{Z}iga Sajovic\\
  \small Polydera\\
  \small \texttt{ziga.sajovic@polydera.com}
}
\date{}


\begin{document}
\maketitle

\begin{abstract}
Mesh CSG pipelines are becoming robust and recently interactive, but unconditional trust is impossible: CSG output crosses the floating-point boundary in both directions, and topology is the only invariant that survives. The pragmatic goal is twofold: guarantee what can be guaranteed, absorb what cannot.

Our pipeline is structured by three layers, each a graph problem on graphs we never explicitly construct. Per-face local exactness --- integer-kernel predicates classify pairwise intersections into five canonical types (VV, VE, VF, EE, EF) without simulation of simplicity; each cut face's 2D arrangement is exact in its own plane. Identity propagation across faces --- all points created along an intersection edge are registered, by edge identity, on every face that contains that edge. Aggregated uncertainty --- when 3D decisions span faces, endpoint materialisation can perturb individual predicates. We exploit the consistency imposed by topological units to absorb local uncertainties: component orientation, relation merge, domain inside/outside.

  The \emph{trueform} implementation delivers robust mesh CSG at interactive speeds in C\texttt{++}, Python, and TypeScript.
\end{abstract}

\section{Introduction}
\label{sec:intro}

Mesh CSG --- arrangements of triangulated surfaces and the booleans that select among them --- underlies geomodeling, CAD, scientific simulation, and graphics. Three decades of work have produced robust pipelines, recently interactive. None is unconditionally trustworthy, and none can be. The pragmatic goal is twofold: guarantee what can be guaranteed, absorb what cannot, by exploiting the consistency topological constraints impose on the noise. Three observations shape the design.

\subsection{Three observations}
\label{sec:observations}

\textbf{Geometric exactness is local.} Predicates must return the exact mathematical sign of a polynomial, not a floating-point answer that depends on the order of summation~\cite{Levy2025}. In a single face's plane, on a bounded integer grid, this is cheap: \texttt{orient3d} and its relatives reduce to fixed-width integer arithmetic with no filter cascade. Several kernels make the same per-call guarantee by different means --- arithmetic expansions~\cite{Shewchuk1997}, indirect predicates~\cite{Attene2020}, multi-precision floats~\cite{Levy2025}, bounded-width integer planes~\cite{Trettner2022}. Extending that exactness from per-call to chained constructions is also possible: materialised intersection points live across face planes, and the arithmetic must follow them --- expansions can reach tens of thousands of components in degenerate cases~\cite{Levy2025}, and fixed-width ladders must widen at each composition or change representation. The cost is real but bounded.

\textbf{The illusion of closed systems.} An arrangement or boolean is never the whole pipeline. Mesh CSG is a step in something larger --- a model on its way to finite-element analysis, mesh editing, or export. Recent systems are exact within their own representation, but use-cases rarely stay within it for more than one operation. Materialising the result back to IEEE double can introduce new intersections that were not in the exact result~\cite{Levy2025}. EMBER acknowledges this in its own limitations, noting that iterated CSG within its representation is exact, but interfacing the floating-point world requires snap rounding which may re-introduce tiny self-intersections~\cite{Trettner2022}. The same boundary works in reverse: inputs come from upstream tools that have already crossed it, so a CSG operation never receives truly clean input either. Pushing precision arbitrarily high inside the system does not change this --- uncertainty re-enters at the boundary regardless of how it was held inside. The pragmatic response is to absorb it within the pipeline, using the structure topology imposes. The only invariant that survives this boundary, in either direction, is topology.

\textbf{Topological constraints absorb uncertainty.} Topology plays two roles. Locally, it is what we guarantee: the contacts the user deliberately specified --- coincident vertices, on-edge points, coplanar faces --- preserved exactly through classification. Globally, it is the structure those contacts implicitly form: manifold-edge connected components, chains of non-manifold edges, paths along which a radial sort must be consistent, domains bounded by selected faces. The local form is what we preserve; the global form is the consistency engine we use to absorb the residual uncertainty the boundary inevitably leaves behind. Treating both as first-class data means accepting non-manifold input as it comes --- flaps, T-junctions, isolated edges --- not as exceptions to be cleaned up before processing. Simulation of Simplicity~\cite{EdelsbrunnerMucke1990} is excluded outright on this ground --- its perturbation strategy is not topologically exact (Figure~\ref{fig:sos_box_divider}). Preserving the local form alone is not enough either. Methods that build wedges edge-by-edge around a non-manifold edge~\cite{BohmRunge2025} preserve local topology yet treat each edge in isolation. The radial sort along such a chain has two failure modes the per-edge view cannot absorb. A single degenerate incident triangle has an ambiguous angular position. And the materialised intersection points the pipeline carries can shift an \texttt{orient3d} wedge enough that a handful of edges in the chain land on the wrong side of a comparison. The bad pairings propagate outward to misclassify valid neighbours. Aggregating observations across the topological unit they belong to --- a component for orientation, a relation for merging, a domain for inside/outside --- is the mitigation. The unit's consistency is the global constraint that pins down the local noise.

\begin{figure}[t]
  \centering
  \def\sclI{4.15}
\def\panelw{7.4}
\def\panelh{5.4}

\definecolor{tfd1face}{HTML}{97F4EB}
\definecolor{tfd1edge}{HTML}{0D665C}
\definecolor{tfd2face}{HTML}{008777}
\definecolor{tfd2edge}{HTML}{063A33}
\definecolor{tfrface}{HTML}{C8467A}
\definecolor{tfredge}{HTML}{5A1F36}
\definecolor{tfvedot}{HTML}{5A1F36}
\definecolor{tfbg}{HTML}{ECECEE}

\newcommand{\proj}[3]{%
  ({((#1) * 0.7880107536 + (#2) * 0.6156614753) * \sclI},
   {((#1) * (-0.1901783216) + (#2) * 0.2433065857 + (#3) * 0.9510565163) * \sclI - 0.35})%
}

\begin{tikzpicture}[font=\sffamily\footnotesize, x=1cm, y=1cm]

  \filldraw[
    fill=tfbg,
    draw=black!20,
    line width=0.4pt,
    rounded corners=3pt
  ]
    (-\panelw/2, -\panelh/2 - 0.35) rectangle (\panelw/2, \panelh/2 - 0.35);

  % --------------------------------------------------------------------
  % D1 outer faces
  % --------------------------------------------------------------------

  \filldraw[
    fill=tfd1face,
    draw=tfd1edge,
    fill opacity=0.40,
    draw opacity=0.75,
    line width=0.45pt
  ]
    \proj{-0.5}{-0.5}{-0.3}
    -- \proj{0}{-0.5}{-0.3}
    -- \proj{0}{0.5}{-0.3}
    -- \proj{-0.5}{0.5}{-0.3}
    -- cycle;

  \filldraw[
    fill=tfd1face,
    draw=tfd1edge,
    fill opacity=0.40,
    draw opacity=0.75,
    line width=0.45pt
  ]
    \proj{-0.5}{-0.5}{0.3}
    -- \proj{0}{-0.5}{0.3}
    -- \proj{0}{0.5}{0.3}
    -- \proj{-0.5}{0.5}{0.3}
    -- cycle;

  \filldraw[
    fill=tfd1face,
    draw=tfd1edge,
    fill opacity=0.40,
    draw opacity=0.75,
    line width=0.45pt
  ]
    \proj{-0.5}{-0.5}{-0.3}
    -- \proj{0}{-0.5}{-0.3}
    -- \proj{0}{-0.5}{0.3}
    -- \proj{-0.5}{-0.5}{0.3}
    -- cycle;

  \filldraw[
    fill=tfd1face,
    draw=tfd1edge,
    fill opacity=0.40,
    draw opacity=0.75,
    line width=0.45pt
  ]
    \proj{0}{0.5}{-0.3}
    -- \proj{-0.5}{0.5}{-0.3}
    -- \proj{-0.5}{0.5}{0.3}
    -- \proj{0}{0.5}{0.3}
    -- cycle;

  \filldraw[
    fill=tfd1face,
    draw=tfd1edge,
    fill opacity=0.40,
    draw opacity=0.75,
    line width=0.45pt
  ]
    \proj{-0.5}{-0.5}{-0.3}
    -- \proj{-0.5}{0.5}{-0.3}
    -- \proj{-0.5}{0.5}{0.3}
    -- \proj{-0.5}{-0.5}{0.3}
    -- cycle;

  % --------------------------------------------------------------------
  % D2 outer faces
  % --------------------------------------------------------------------

  \filldraw[
    fill=tfd2face,
    draw=tfd2edge,
    fill opacity=0.40,
    draw opacity=0.75,
    line width=0.45pt
  ]
    \proj{0}{-0.5}{-0.3}
    -- \proj{0.5}{-0.5}{-0.3}
    -- \proj{0.5}{0.5}{-0.3}
    -- \proj{0}{0.5}{-0.3}
    -- cycle;

  \filldraw[
    fill=tfd2face,
    draw=tfd2edge,
    fill opacity=0.40,
    draw opacity=0.75,
    line width=0.45pt
  ]
    \proj{0}{-0.5}{0.3}
    -- \proj{0.5}{-0.5}{0.3}
    -- \proj{0.5}{0.5}{0.3}
    -- \proj{0}{0.5}{0.3}
    -- cycle;

  \filldraw[
    fill=tfd2face,
    draw=tfd2edge,
    fill opacity=0.40,
    draw opacity=0.75,
    line width=0.45pt
  ]
    \proj{0}{-0.5}{-0.3}
    -- \proj{0.5}{-0.5}{-0.3}
    -- \proj{0.5}{-0.5}{0.3}
    -- \proj{0}{-0.5}{0.3}
    -- cycle;

  \filldraw[
    fill=tfd2face,
    draw=tfd2edge,
    fill opacity=0.40,
    draw opacity=0.75,
    line width=0.45pt
  ]
    \proj{0.5}{0.5}{-0.3}
    -- \proj{0}{0.5}{-0.3}
    -- \proj{0}{0.5}{0.3}
    -- \proj{0.5}{0.5}{0.3}
    -- cycle;

  \filldraw[
    fill=tfd2face,
    draw=tfd2edge,
    fill opacity=0.40,
    draw opacity=0.75,
    line width=0.45pt
  ]
    \proj{0.5}{-0.5}{-0.3}
    -- \proj{0.5}{0.5}{-0.3}
    -- \proj{0.5}{0.5}{0.3}
    -- \proj{0.5}{-0.5}{0.3}
    -- cycle;

  % --------------------------------------------------------------------
  % Cutting rectangle R
  % --------------------------------------------------------------------

  \filldraw[
    fill=tfrface,
    draw=tfredge,
    fill opacity=0.70,
    draw opacity=0.95,
    line width=0.9pt
  ]
    \proj{0}{-0.5}{-0.3}
    -- \proj{0}{0.5}{-0.3}
    -- \proj{0}{0.5}{0.3}
    -- \proj{0}{-0.5}{0.3}
    -- cycle;

  % --------------------------------------------------------------------
  % VE intersection points
  % --------------------------------------------------------------------

  \foreach \x/\y/\z in {
    0/-0.5/-0.3,
    0/0.5/-0.3,
    0/0.5/0.3,
    0/-0.5/0.3
  }{
    \filldraw[
      fill=tfvedot,
      draw=white,
      line width=0.45pt
    ] \proj{\x}{\y}{\z} circle [radius=3.2pt];
  }

  % --------------------------------------------------------------------
  % Legend above drawing
  % --------------------------------------------------------------------

  \node[
    anchor=south,
    fill=tfbg,
    draw=black!18,
    rounded corners=1pt,
    inner xsep=5pt,
    inner ysep=4pt,
    font=\footnotesize,
    minimum width=\panelw cm
  ] at (0, 2.55)
  {
    \begin{tabular}{@{}l@{\hspace{0.4em}}l@{\hspace{1.6cm}}l@{\hspace{0.4em}}l@{}}
      \tikz{\draw[fill=tfd1face, draw=tfd1edge, fill opacity=0.40] (0,0) rectangle (0.28,0.16);}
      &
      Sub-domain $D_1$
      &
      \tikz{\draw[fill=tfd2face, draw=tfd2edge, fill opacity=0.40] (0,0) rectangle (0.28,0.16);}
      &
      Sub-domain $D_2$
      \\[2pt]
      \tikz{\draw[fill=tfrface, draw=tfredge, fill opacity=0.70] (0,0) rectangle (0.28,0.16);}
      &
      Cutting rectangle $R$
      &
      \tikz{\filldraw[fill=tfvedot, draw=white, line width=0.4pt] (0.14,0.08) circle (0.10);}
      &
      VE intersection
    \end{tabular}
  };

\end{tikzpicture}

  \caption{Mesh arrangements often divide a volume into regions of interest. When those regions are cut parametrically, primitives end up coplanar by construction --- as here, where a box is split into sub-domains $D_1, D_2$ by an internal rectangle $R$ whose vertices sit at midpoints of cube edges (four VE intersections). Simulation of simplicity resolves these coplanar contacts arbitrarily, and therefore does not reliably produce correct domains.}
  \label{fig:sos_box_divider}
\end{figure}

\subsection{Contributions}
\label{sec:contributions}


\begin{figure*}[t]
  \centering
  \def\sclI{0.7}
\def\sclF{0.9}
\def\panelw{5.2}
\def\panelh{4.4}
\def\panelgap{0.55}

\newcommand{\iso}[3]{({((#1) - (#2)) * 0.866 * \sclI}, {((#3) - ((#1) + (#2)) * 0.5) * \sclI})}

\begin{tikzpicture}[font=\sffamily\footnotesize, x=1cm, y=1cm]

  % =====================================================================
  % PANEL 1 -- 3D configuration
  % =====================================================================
  \begin{scope}[shift={(0,0)}]
    \filldraw[fill=black!3, draw=black!20, line width=0.4pt, rounded corners=3pt]
      (-\panelw/2, -\panelh/2) rectangle (\panelw/2, \panelh/2);

    % Quad D (back, dimmest)
    \fill[tfteal, opacity=0.10]
      \iso{-2}{0.3}{-1.3} -- \iso{2}{0.3}{-1.3} -- \iso{2}{0.3}{1.3} -- \iso{-2}{0.3}{1.3} -- cycle;
    \draw[tfteal, line width=0.5pt, opacity=0.55]
      \iso{-2}{0.3}{-1.3} -- \iso{2}{0.3}{-1.3} -- \iso{2}{0.3}{1.3} -- \iso{-2}{0.3}{1.3} -- cycle;

    % Quad B
    \fill[tfteal, opacity=0.13]
      \iso{-0.8}{-2}{-1.3} -- \iso{-0.8}{2}{-1.3} -- \iso{-0.8}{2}{1.3} -- \iso{-0.8}{-2}{1.3} -- cycle;
    \draw[tfteal, line width=0.5pt, opacity=0.65]
      \iso{-0.8}{-2}{-1.3} -- \iso{-0.8}{2}{-1.3} -- \iso{-0.8}{2}{1.3} -- \iso{-0.8}{-2}{1.3} -- cycle;

    % Quad C
    \fill[tfteal, opacity=0.13]
      \iso{0.8}{-2}{-1.3} -- \iso{0.8}{2}{-1.3} -- \iso{0.8}{2}{1.3} -- \iso{0.8}{-2}{1.3} -- cycle;
    \draw[tfteal, line width=0.5pt, opacity=0.65]
      \iso{0.8}{-2}{-1.3} -- \iso{0.8}{2}{-1.3} -- \iso{0.8}{2}{1.3} -- \iso{0.8}{-2}{1.3} -- cycle;

    % Quad A (front, most opaque)
    \fill[tfteal, opacity=0.20]
      \iso{-2}{-2}{0} -- \iso{-2}{2}{0} -- \iso{2}{2}{0} -- \iso{2}{-2}{0} -- cycle;
    \draw[tfteal, line width=0.7pt, opacity=0.85]
      \iso{-2}{-2}{0} -- \iso{-2}{2}{0} -- \iso{2}{2}{0} -- \iso{2}{-2}{0} -- cycle;

    % Intersection edges on A
    \draw[tfrose, line width=1pt, opacity=0.95] \iso{-0.8}{-2}{0} -- \iso{-0.8}{2}{0};
    \draw[tfrose, line width=1pt, opacity=0.95] \iso{0.8}{-2}{0}  -- \iso{0.8}{2}{0};
    \draw[tfrose, line width=1pt, opacity=0.85] \iso{-2}{0.3}{0}  -- \iso{2}{0.3}{0};

    % Boundary vertices
    \fill[tfteal] \iso{-0.8}{-2}{0} circle [radius=1.4pt];
    \fill[tfteal] \iso{-0.8}{2}{0}  circle [radius=1.4pt];
    \fill[tfteal] \iso{0.8}{-2}{0}  circle [radius=1.4pt];
    \fill[tfteal] \iso{0.8}{2}{0}   circle [radius=1.4pt];
    \fill[tfteal] \iso{-2}{0.3}{0}  circle [radius=1.4pt];
    \fill[tfteal] \iso{2}{0.3}{0}   circle [radius=1.4pt];

    % Quad labels — anchored at the front/top corners of each quad
    \node[text=black!75, font=\sffamily\small\bfseries, anchor=south east]
      at \iso{2.1}{-2}{0.3} {$A$};
    \node[text=black!55, font=\sffamily\small\bfseries\itshape, anchor=north west]
      at \iso{-0.8}{-2}{1.3} {$B$};
    \node[text=black!55, font=\sffamily\small\bfseries\itshape, anchor=south west]
      at \iso{0.8}{-2}{1.3} {$C$};
    \node[text=black!55, font=\sffamily\small\bfseries\itshape, anchor=south east]
      at \iso{-2}{0.3}{1.3} {$D$};

  \end{scope}

  % =====================================================================
  % PANEL 2 -- edges on face A
  % =====================================================================
  \begin{scope}[shift={(\panelw + \panelgap, 0)}]
    \filldraw[fill=black!3, draw=black!20, line width=0.4pt, rounded corners=3pt]
      (-\panelw/2, -\panelh/2) rectangle (\panelw/2, \panelh/2);

    % Face A as a square (scaled to sclF)
    \fill[tfteal, opacity=0.08] (-2*\sclF, -2*\sclF) rectangle (2*\sclF, 2*\sclF);
    \draw[tfteal, line width=0.7pt, opacity=0.75] (-2*\sclF, -2*\sclF) rectangle (2*\sclF, 2*\sclF);

    % 3 edges
    \draw[tfrose, line width=1pt, opacity=0.95] (-0.8*\sclF, -2*\sclF) -- (-0.8*\sclF, 2*\sclF);
    \draw[tfrose, line width=1pt, opacity=0.95] ( 0.8*\sclF, -2*\sclF) -- ( 0.8*\sclF, 2*\sclF);
    \draw[tfrose, line width=1pt, opacity=0.85] (-2*\sclF,  0.3*\sclF) -- (2*\sclF,  0.3*\sclF);

    % 6 boundary vertices
    \fill[tfteal] (-0.8*\sclF, -2*\sclF)  circle [radius=1.6pt];
    \fill[tfteal] (-0.8*\sclF,  2*\sclF)  circle [radius=1.6pt];
    \fill[tfteal] ( 0.8*\sclF, -2*\sclF)  circle [radius=1.6pt];
    \fill[tfteal] ( 0.8*\sclF,  2*\sclF)  circle [radius=1.6pt];
    \fill[tfteal] (-2*\sclF,    0.3*\sclF) circle [radius=1.6pt];
    \fill[tfteal] ( 2*\sclF,    0.3*\sclF) circle [radius=1.6pt];

    % Edge labels (black/dark) — placed by the edges they mark
    \node[text=black!70, font=\sffamily\scriptsize\ttfamily]
      at (-\sclF - 0.4, \sclF) {(A,B)};
    \node[text=black!70, font=\sffamily\scriptsize\ttfamily]
      at ( \sclF + 0.4, \sclF) {(A,C)};
    \node[text=black!70, font=\sffamily\scriptsize\ttfamily, anchor=south]
      at (0, 0.3*\sclF + 0.08) {(A,D)};
  \end{scope}

  % =====================================================================
  % PANEL 3 -- intersection graph
  % =====================================================================
  \begin{scope}[shift={(2*\panelw + 2*\panelgap, 0)}]
    \filldraw[fill=black!3, draw=black!20, line width=0.4pt, rounded corners=3pt]
      (-\panelw/2, -\panelh/2) rectangle (\panelw/2, \panelh/2);

    % Faded face boundary
    \fill[tfteal, opacity=0.04] (-2*\sclF, -2*\sclF) rectangle (2*\sclF, 2*\sclF);
    \draw[tfteal, line width=0.4pt, opacity=0.30] (-2*\sclF, -2*\sclF) rectangle (2*\sclF, 2*\sclF);

    % B: 2 segments (split at y=0.3)
    \draw[tfrose, line width=0.9pt, opacity=0.70] (-0.8*\sclF, -2*\sclF)  -- (-0.8*\sclF, 0.3*\sclF);
    \draw[tfrose, line width=0.9pt, opacity=0.70] (-0.8*\sclF,  0.3*\sclF) -- (-0.8*\sclF, 2*\sclF);

    % C: 2 segments
    \draw[tfrose, line width=0.9pt, opacity=0.70] ( 0.8*\sclF, -2*\sclF)  -- ( 0.8*\sclF, 0.3*\sclF);
    \draw[tfrose, line width=0.9pt, opacity=0.70] ( 0.8*\sclF,  0.3*\sclF) -- ( 0.8*\sclF, 2*\sclF);

    % D: 3 segments (split at x=-0.8 and x=0.8)
    \draw[tfrose, line width=0.9pt, opacity=0.60] (-2*\sclF,    0.3*\sclF) -- (-0.8*\sclF, 0.3*\sclF);
    \draw[tfrose, line width=0.9pt, opacity=0.60] (-0.8*\sclF,  0.3*\sclF) -- ( 0.8*\sclF, 0.3*\sclF);
    \draw[tfrose, line width=0.9pt, opacity=0.60] ( 0.8*\sclF,  0.3*\sclF) -- ( 2*\sclF,   0.3*\sclF);

    % Endpoints (faded teal)
    \fill[tfteal, opacity=0.50] (-0.8*\sclF, -2*\sclF)  circle [radius=1.1pt];
    \fill[tfteal, opacity=0.50] (-0.8*\sclF,  2*\sclF)  circle [radius=1.1pt];
    \fill[tfteal, opacity=0.50] ( 0.8*\sclF, -2*\sclF)  circle [radius=1.1pt];
    \fill[tfteal, opacity=0.50] ( 0.8*\sclF,  2*\sclF)  circle [radius=1.1pt];
    \fill[tfteal, opacity=0.50] (-2*\sclF,    0.3*\sclF) circle [radius=1.1pt];
    \fill[tfteal, opacity=0.50] ( 2*\sclF,    0.3*\sclF) circle [radius=1.1pt];

    % Crossings (rose, prominent)
    \fill[tfrose] (-0.8*\sclF, 0.3*\sclF) circle [radius=2pt];
    \fill[tfrose] ( 0.8*\sclF, 0.3*\sclF) circle [radius=2pt];

    % Triplet labels (black, near each crossing)
    \node[text=black!75, font=\sffamily\scriptsize\ttfamily\bfseries, anchor=south east]
      at (-0.8*\sclF - 0.04, 0.3*\sclF + 0.05) {\{A,B,D\}};
    \node[text=black!75, font=\sffamily\scriptsize\ttfamily\bfseries, anchor=south west]
      at ( 0.8*\sclF + 0.04, 0.3*\sclF + 0.05) {\{A,C,D\}};

  \end{scope}

\end{tikzpicture}

  \caption{Two-stage local arrangement on face $A$, which meets faces $B$, $C$, $D$. \emph{Left}: 3D configuration with three intersection edges and six boundary vertices on $A$. \emph{Middle}: those edges in $A$'s local plane, each named by the face pair that produced it (\texttt{(A,B)}, \texttt{(A,C)}, \texttt{(A,D)}). \emph{Right}: the planar arrangement --- seven segments, two crossings, each crossing named by its face-pair triplet (\texttt{\{A,B,D\}}, \texttt{\{A,C,D\}}); points computed independently from different face slices unify without coordinate comparison.}
  \label{fig:arrangement_pipeline}
\end{figure*}

We make five contributions:

\begin{enumerate}
\item \textbf{Topologically exact predicates over a bounded integer kernel.}
Local topological intent --- the coincident vertices, on-edge points, and coplanar contacts the user specified --- is preserved exactly through five canonical intersection types (VV, VE, VF, EE, EF), classified by exact integer predicates in a single bounded integer kernel (Section~\ref{sec:predicates}). Classification is decided by the zero-pattern of \texttt{orient3d} signs across a fan triangulation of the polygon, distinguishing real polygon edges from internal fan diagonals, \emph{before} any point is constructed. The five-type discipline preserves what Simulation of Simplicity~\cite{EdelsbrunnerMucke1990} discards.

\item \textbf{A two-stage local arrangement.}
The arrangement is built without constructing a global structure (Section~\ref{sec:arrangement}). Stage one computes pairwise intersections in parallel, naming points by their originating face-pair triplets and intersection edges by their two endpoints --- an identity that is face-independent. Stage two computes an exact 2D segment-arrangement inside each face's local plane (Figure~\ref{fig:arrangement_pipeline}); topological exactness of the VV, VE, and EE classification ensures point and segment uniqueness. The two stages are coupled only through this identity: constructed points computed independently from different face slices unify; any new point materialised on a segment between two named endpoints is registered on every face containing it. The same code path handles two-mesh, $N$-mesh, and polygon-soup self-arrangement with no special casing. Per-face $n$ stays small in practice, bounded by local cut geometry; what scales in real workloads is the number of cut faces, parallel across threads.

\item \textbf{Implicit graph structures.}
Every stage of the pipeline is a graph problem on graphs that are never explicitly constructed (Section~\ref{sec:graphs}). Flat array layouts (offsets, dense ranks, adjacency tables, equivalence maps) indexed by lookup carry the global consistency constraints topology imposes. Nodes of the reduced graph are manifold-edge connected components of intersected and non-intersected regions of the input mesh; coplanar contact regions become a single component by construction, absorbing the separate scar-removal pass L\'evy performs~\cite{Levy2025}. Non-manifold flaps, T-junctions, and isolated edges flow through first-class. The reduced graph carries the combinatorial information of L\'evy's Weiler 3-map~\cite{Weiler1985} with no half-edge structure to maintain. L\'evy~\cite{Levy2025} advocates exploiting the combinatorics; we sharpen the principle by treating every stage as a graph problem on implicit graphs.

\item \textbf{Aggregated uncertainty across topological units.}
When 3D decisions span faces, the materialised intersection points the pipeline carries can perturb individual \texttt{orient3d} wedge comparisons. Aggregation absorbs the local noise at each topological unit (Section~\ref{sec:aggregation}): components vote on orientation, relations between components vote on merging into domains, and domains vote on inside/outside. All three are weighted-majority votes over exact \texttt{orient3d} observations, sidestepping the per-patch tangency handling that ray-casting requires~\cite{Cherchi2022}. Ambiguous local observations are skipped; the faces involved still belong to their component topologically and inherit the label. Domains carrying no seeds fall back to ray-casting.

\item \textbf{A library for interactive exact CSG in C\texttt{++}, Python, and the browser.}
A single header-only C\texttt{++} library exposes the pipeline through Python (nanobind) and TypeScript (WebAssembly) bindings (Section~\ref{sec:implementation}). The same exact kernel runs on x86, ARM, and in the browser, at interactive speeds. The codebase is freely available for academic and research use~\cite{trueform2025}. To our knowledge, no prior CSG library is interactive in the browser.
\end{enumerate}

\subsection{Related work}
\label{sec:related}

\textbf{Geometric exactness.} The arithmetic-kernel literature has converged on a small set of approaches. Shewchuk's adaptive floating-point expansions~\cite{Shewchuk1997} compute a predicate's sign exactly by accumulating non-overlapping FP components, refining only when the filter cascade fails. FPG~\cite{MeyerPion2008} generates such filtered predicates from polynomial expressions. CGAL's lazy rationals~\cite{CGAL2023} offer the same exactness through arbitrary-precision evaluation. Indirect predicates~\cite{Attene2020} take a different route: they store constructed points as recipes (line-plane intersections, triple-plane intersections) and fold them into the predicate's polynomial. The FA mesh arrangement~\cite{Cherchi2020} is built on this. EMBER's plane-based representation~\cite{Trettner2022,NehringWirxel2021} uses fixed-width 256-bit integer arithmetic with no filter cascade. L\'evy's recent CSG pipeline~\cite{Levy2025} carries two kernels in parallel: Shewchuk expansions for the easy cases, GMP multiprecision for the hard ones. Interval arithmetic~\cite{Bronnimann1998} provides the cheap filter all of these depend on. Trueform sits between Attene and EMBER: a fixed three-width integer ladder (T0$\to$T1$\to$T2), no dynamic allocation, no fallback kernel.

\textbf{Open systems and the float boundary.} Once exact coordinates materialise to floating point, new intersections may appear. 3D snap rounding~\cite{DevillersLazardLenhart2018} formalises the problem and gives an algorithm; Valque~\cite{Valque2019} pushes it toward practical implementation. Geometric rounding for meshes generalises the same problem class~\cite{MilenkovicSacks2019}. Cherchi's mesh arrangement applies an iterated cast-to-float heuristic to mitigate the damage~\cite{Cherchi2020}. EMBER explicitly acknowledges that interfacing the floating-point world reintroduces tiny self-intersections~\cite{Trettner2022}. None of these approaches claim to remove the boundary --- they minimise its damage. Trueform accepts the boundary as given and routes uncertainty into topological aggregation.

\textbf{Topology.} Mesh arrangements compute the global structure that classification consumes. Zhou's pipeline~\cite{Zhou2016} uses CGAL exact predicates and topological flooding. Cherchi 2020~\cite{Cherchi2020} replaces this with indirect predicates and ear-cut triangulation; the interactive 2022 variant~\cite{Cherchi2022} keeps the arrangement and uses a per-patch ray cast. L\'evy~\cite{Levy2025} builds the Weiler 3-map~\cite{Weiler1985} in the language of combinatorial maps~\cite{Lienhardt1988}, then floods bit-vectors of operand membership through it. Bohm and Runge~\cite{BohmRunge2025} extract domains from non-manifold meshes by per-edge angular sort and wedge construction --- a 3D generalisation of plane-graph region extraction~\cite{JiangBunke1993}. Generalized winding numbers~\cite{JacobsonKavanSorkine2013} provide an orthogonal path to inside/outside that tolerates non-manifold input. Simulation of Simplicity~\cite{EdelsbrunnerMucke1990} sidesteps the global construction by symbolic perturbation; it is the topology-discarding baseline of \S\ref{sec:observations}.

\textbf{Application context.} Mesh CSG drives demanding workflows in geomodelling~\cite{Caumon2003,Pellerin2017}, interactive mechanical modelling~\cite{Tymms2016}, and CSG for fabrication~\cite{Garg2016}. All share the same constraint from \S\ref{sec:observations}: CSG output is consumed by another stage in floating point, and that stage cannot tolerate topology drift.

\subsection{Algorithm overview}
\label{sec:overview}

The pipeline is a single forward pass. Stage one intersects every candidate face pair and emits per-face intersection records (Section~\ref{sec:predicates},~\ref{sec:arrangement}). Stage two arranges each cut face's records independently in its local plane; closed inner loops are patched into the regions that contain them (Section~\ref{sec:arrangement}). Edge adjacency is built only for the new cut sub-faces; non-intersected original faces reuse the input mesh's manifold-edge link directly. Together they form the implicit reduced graph that classification operates on.

Classification proceeds in three aggregation steps (Section~\ref{sec:aggregation}). First, each component aggregates orientation-flip votes across its faces. Second, each pair of components related by a chain of non-manifold edges aggregates wedge merge votes along the chain; union-find then propagates the merge decisions on the reduced graph to yield volumetric domains. Third, each domain aggregates wedge inside/outside seeds drawn along its non-manifold edge paths. All three are weighted-majority votes over exact \texttt{orient3d} observations; ambiguous observations are skipped, with faces inheriting their component's label. Domains carrying no seeds fall back to ray-casting.

The boolean operator selects domains by their classifications. The final mesh is assembled only at the end: selected uncut original faces pass through directly, and selected cut sub-faces are triangulated into the output.


\section{The algorithm}
\label{sec:algorithm}

\subsection{Topologically exact predicates on convex polygons}
\label{sec:predicates}

The kernel is a bounded integer ladder over snapped coordinates. Input floats are mapped onto an integer grid. Predicate expressions promote width as needed: subtractions cost one bit (T0$\to$T1), products of differences double width (T1$\times$T1$\to$T2). The ladder closes after T2. Every predicate the pipeline uses --- \texttt{orient2d}, \texttt{orient3d}, and the cross/dot work that feeds them --- is a polynomial of total degree at most three. The library exposes two ladders, $32{\to}64{\to}128$ and $64{\to}128{\to}256$ bits; the user picks based on input precision and resolution. The kernel is the substrate; Attene's indirect predicates~\cite{Attene2020} carry the same accounting via stored expressions, EMBER~\cite{Trettner2022} via fixed-width 256-bit planes. Trueform works in between, with a fixed promotion ladder of exactly known widths and no dynamic allocation.

Five canonical configurations carry the local topology forward: vertex-vertex (\texttt{VV}), vertex-edge (\texttt{VE}), vertex-face (\texttt{VF}), edge-edge (\texttt{EE}), and edge-face (\texttt{EF}). Each emitted intersection is exactly one of these five, named by the simplex pair that produced it. The same taxonomy lives in L\'evy~\cite{Levy2025} as $(\sigma, \sigma')$ pairs over a triangle's seven open simplices; Trueform generalises to convex polygons.

Convex polygons are the input scope because they fan-triangulate trivially online from any vertex $A_0$. There is no separate CDT, no auxiliary index, no triangulation queue. The fan triangulation is what makes the rest of the section possible. Input polygons cannot be assumed to be in general position; three or more colinear vertices occur in practice. A fan triangle is the right unit of work. A non-degenerate triangle classifies cleanly via the zero-pattern below; a degenerate one (three colinear vertices) contributes nothing to that slice, and the loop advances. Non-convex face primitives are rare in mesh CSG inputs, so the convex scope captures the practical case at minimal cost.

The classifier evaluates three \texttt{orient3d} signs per fan triangle. For fan triangle $t = (A_0, A_{t+1}, A_{t+2})$ and segment $(D, E)$, the signs are
\[
\begin{aligned}
v_1 &= \mathrm{orient3d}(A_0,     A_{t+1}, D, E),\\
v_2 &= \mathrm{orient3d}(A_{t+1}, A_{t+2}, D, E),\\
v_3 &= \mathrm{orient3d}(A_0,     A_{t+2}, D, E).
\end{aligned}
\]
The zero-pattern decides the type. No zeros and a sign disagreement on the polygon edge: \texttt{EF}. One zero on a real polygon edge: \texttt{VE} (or \texttt{EE} when both segment ends sit on the same edge). Two zeros at a shared fan vertex: \texttt{VV}. Fan diagonals --- the edges $A_0 A_{t+1}$ for $t > 0$ and $A_0 A_{t+2}$ for $t + 3 < n$ --- are not real polygon edges. A zero on one of them is geometrically real but topologically a fan-interior point, classified as \texttt{EF}, never as \texttt{VE}. The real-edge predicate $(v_1 = 0 \land t = 0) \lor (v_2 = 0) \lor (v_3 = 0 \land t + 3 = n)$ is the only place the fan triangulation leaks structural information; the rest is rotation-symmetric.

Two complementary routines run over this decomposition. The first processes edges that strictly cross the opposing plane, emitting \texttt{EF}, \texttt{EE}, or \texttt{VE} from the fan zero-pattern above. The second processes on-plane primitives --- any vertex whose plane-side sign is zero --- and emits \texttt{VV}, coplanar \texttt{VE}, or coplanar \texttt{EE}. A third routine emits \texttt{VF}. The three are not mutually exclusive: for a pair where some vertices cross and some touch the plane, all three may fire. The partition into types is still exact --- every real contact is emitted once, from one routine.

Degenerate polygons are not a special case. A chain of colinear vertices is a polygon that has shrunk to an edge. When the chain lies in the opposing plane, every vertex reads sign zero. The on-plane routine then catches all its contacts as the same \texttt{VV}/\texttt{VE}/\texttt{EE} primitives a real edge would produce. Convexity makes the geometry simple; the classifier makes the topology uniform. Simulation of Simplicity~\cite{EdelsbrunnerMucke1990} is not needed at any step.

\subsection{Two-stage local arrangement}
\label{sec:arrangement}

The arrangement is computed in two stages. Stage one intersects every candidate face pair in parallel, emitting tagged intersection records keyed by simplex pair on each face. Stage two arranges each cut face's records independently in its local 2D plane. The implementation exposes the arrangement either over a collection of meshes or over a polygon soup; the algorithm is the same, only the candidate set differs.

\subsubsection*{Stage one --- pairwise intersection records}

Candidate face pairs come from a parallel AABB-tree descent --- dual across two trees for the multi-mesh case, self over a single tree for polygon soup. Each surviving pair invokes the predicates of Section~\ref{sec:predicates}. The output is a flat tuple per intersection: which mesh, which polygon, which polygon on the other side, the simplex pair on each, the point identity. No edges, no per-face geometry yet --- just identities.

The stream is globally unique without any dedup pass. Two mechanisms ensure this. \emph{Within a pair}, the decision tree of Section~\ref{sec:predicates} fires exactly one type per real contact --- no overlapping sub-simplex queries like L\'evy's~\cite{Levy2025}, which require a \texttt{sort + unique} pass. \emph{Across pairs}, each vertex and edge of a mesh has one designated owner face, picked up front from the face-membership list and the manifold-edge link. During traversal, only the owner emits records for that primitive; the other faces that share it skip it. Together the two mechanisms keep every contact recorded exactly once, regardless of how many adjacent faces share the primitive.

Construction is decoupled from classification. \texttt{VV}, \texttt{VE}, and \texttt{VF} reuse input-vertex identities directly. Only \texttt{EE} and \texttt{EF} compute a new point, via barycentric-weighted \texttt{div\_round} on the integer endpoints. The new point lives at exact integer coordinates on the T2 grid. Its identity is the canonical triplet of input faces that produced it --- a global key independent of where it was computed.

A parallel sort by key $(\text{tag}, \text{polygon}, \text{tag}', \text{polygon}')$ groups records. Offset-blocks on the prefix $(\text{tag}, \text{polygon})$ place every record in a contiguous slice belonging to one face. Stage two runs over these slices in parallel.

\subsubsection*{Stage two --- per-face planar arrangement}

Per-face processing has three phases: segment extraction, segment arrangements, and region extraction. The first two build a per-face segment graph; the third walks regions out of it.

All work happens in the local 2D plane of each face. Every intersection point --- stage-one outputs and the new crossings stage two produces --- carries a canonical identity and a 2D copy in each face that contains it. These per-face 2D copies let per-face arithmetic stay exact: coplanarity holds within each plane by construction, and the integer planar predicates of \S\ref{sec:predicates} apply without further constraint. Materialisation to a 3D output happens once, at the end of aggregation, where the open-system limitation of \S\ref{sec:observations} bites at most once.

\paragraph{Segment extraction.}
Extraction proceeds in two passes. The first inserts the face's intersection points into its base loop, ordered around the perimeter. The result is the cut base loop. The second walks the face's slice. The stage-one sort ordered records by $(\text{tag}, \text{polygon}, \text{tag}', \text{polygon}')$, so records sharing a cross-face context are already contiguous --- the grouping is the linear sequence itself. Each run classifies in $O(1)$ from its size: one record means no edge; two means a single intersection edge; three or more mean a closed coplanar polygon-of-contact, which is a contiguous subset of the other polygon's cut base loop. The subset is traced linearly through the triplet identities of its points; no geometric check is needed. Each intersection edge has a canonical identity given by its two endpoint identities. The same canonical edge appears in every face that contains it; updates to the edge propagate to all containing faces by lookup.

\paragraph{Segment arrangements.}
A per-face planar arrangement resolves crossings among the extracted segments. Each crossing produces a new intersection point, identified by the canonical intersection edges it lies on.

Faces with fewer than 32 segments run a quadratic crossing check; larger faces build a segment tree and find crossings in $O(n \log n)$. Each new crossing is classified by three types (VV, VE, EE), the 2D restriction of the five-type discipline of stage one. Topological exactness of the three types ensures point and segment uniqueness: coincident endpoints --- points landing on the same integer-grid cell --- collapse into a single point as VV, and overlapping segments collapse into one canonical segment. The records on any single face are typically a handful, occasionally a few dozen --- bounded by the local cut geometry --- so $n$ is small. What scales in real workloads is the number of cut faces, parallelised across threads. The strategy is deliberate: keep the heavy algorithmic step inside a small per-face set, and ride throughput on face-level parallelism rather than per-face triangulation.

Identity stays globally consistent through a canonicalisation step that sorts every edge instance by canonical endpoint pair and groups all per-face copies of the same intersection edge into a single offset block --- the reverse map from canonical edges to containing faces, built by structure alone. Each canonical intersection edge is then processed once, in parallel across blocks, to absorb the crossings naming it: crossings are sorted by parametric position $t$ along the segment, the segment is split into sub-segments at the sorted positions, and two crossings landing at the same $t$ merge into a single point by the same VV mechanism. The split rebuilds the shared buffer in place; every face's slice then sees the new sub-segments by lookup. Point identities are remapped in a single parallel pass over the buffer, so the next phase sees a globally consistent segment graph on every face.

\paragraph{Region extraction.}
With the segment graph in place, regions of the cut face follow from a path classification. Segments form paths in the graph. A path is a \emph{crossing} when both endpoints sit on the cut base loop. It is \emph{non-crossing} when its terminus is an internal junction of degree three or more --- a condition that arises only at points where three or more faces meet, or at non-manifold edge junctions. The majority of faces produce only crossing paths.

Region extraction dispatches on this distinction. When every path is a crossing, regions are extracted in linear time by dividing the cut base loop at each crossing into sub-loops. The trivial common case --- a single crossing edge --- yields two sub-loops with no predicate work at all. When non-crossing paths are present, an exact planar walk runs on the full edge set: half-edges at each vertex are placed in radial order by an exact \texttt{orient2d} comparator, and the regions are walked.

Closed paths and dangling cut paths are patched into the regions that contain them by an exact point-in-region test. The output of stage two is the set of sub-loops belonging to each cut face.

\subsection{Implicit graph structures}
\label{sec:graphs}

\subsubsection*{Combinatorics over geometry, at every stage}

\subsubsection*{The reduced graph: manifold-edge components as nodes}

\subsubsection*{Exact radial sort and region walk}

\subsubsection*{Combinatorial equivalence to the Weiler 3-map}

\subsection{Aggregation across topological units}
\label{sec:aggregation}

\subsubsection*{The three aggregation steps}

\subsubsection*{Component orientation by weighted-majority flip}

\subsubsection*{Domain extraction by radial sort and union-find}

\subsubsection*{Per-domain classification by wedge seeds}

\subsubsection*{Skip-and-inherit; ray-cast fallback for orphan domains}

\subsection{Implementation}
\label{sec:implementation}

\section{Results}
\label{sec:results}

\section{Conclusions}
\label{sec:conclusions}

\section*{Acknowledgements}

\bibliographystyle{plainnat}
\bibliography{references}

\end{document}
